\documentclass[UTF8,12pt]{ctexart}
\usepackage[a4paper,margin=24mm]{geometry}
\usepackage{amsmath,amssymb,bm,graphicx,adjustbox}
\newcommand{\fitmath}[1]{\begin{adjustbox}{max width=\linewidth}$\displaystyle #1$\end{adjustbox}}
\setlength{\parindent}{0pt}
\setlength{\parskip}{.7em}
\sloppy
\allowdisplaybreaks
\begin{document}
\section*{2008 年考研数学三 · 参考答案}
\subsection*{2008年真题参考答案}

\subsection*{一、选择题}

（1）B。（2）C。（3）B。（4）A。（5）C。（6）D。（7）A。（8）D。


\subsection*{二、填空题}

（9）\(\displaystyle 1\)。（10）\(\displaystyle \dfrac{\displaystyle 1}{\displaystyle 2}\ln3\)。（11）\(\displaystyle \dfrac{\displaystyle \pi}{\displaystyle 4}\)。（12）\(\displaystyle \dfrac{\displaystyle 1}{\displaystyle x}\)。（13）\(\displaystyle 3\)。（14）\(\displaystyle \dfrac{\displaystyle 1}{\displaystyle 2e}\)。


\subsection*{三、解答题}

（15）\(\displaystyle -\dfrac{\displaystyle 1}{\displaystyle 6}\)。


（16）（I）\(\displaystyle \mathrm dz=\dfrac{\displaystyle 2x-\varphi^{\prime}(x+y+z)}{\displaystyle \varphi^{\prime}(x+y+z)+1}\,\mathrm dx+\dfrac{\displaystyle 2y-\varphi^{\prime}(x+y+z)}{\displaystyle \varphi^{\prime}(x+y+z)+1}\,\mathrm dy\)。（II）\(\displaystyle \dfrac{\displaystyle \partial u}{\displaystyle \partial x}=\dfrac{\displaystyle -2(2x+1)\varphi^{\prime\prime}(x+y+z)}{\displaystyle [\varphi^{\prime}(x+y+z)+1]^3}\)。


（17）\(\displaystyle \dfrac{\displaystyle 19}{\displaystyle 4}+\ln2\)。（18）证明略。（19）\(\displaystyle A\) 至少应为3980万元。


（20）（I）证明略。（II）\(\displaystyle a\ne0\) 时，方程组有唯一解，\(\displaystyle x_1=\dfrac{\displaystyle n}{\displaystyle (n+1)a}\)。（III）\(\displaystyle a=0\) 时，方程组有无穷多解，通解为 \(\displaystyle x=k(1,\allowbreak 0,\allowbreak \cdots,\allowbreak 0)^{\mathrm T}+(0,\allowbreak 1,\allowbreak 0,\allowbreak \cdots,\allowbreak 0)^{\mathrm T}\)，其中 \(\displaystyle k\) 为任意常数。


（21）（I）证明略。（II）\(\displaystyle P^{-1}AP=\begin{pmatrix}-1&0&0\\0&1&1\\0&0&1\end{pmatrix}\)。


（22）（I）\(\displaystyle \dfrac{\displaystyle 1}{\displaystyle 2}\)。（II）\(\displaystyle f_Z(z)=\begin{cases}\dfrac{\displaystyle 1}{\displaystyle 3},\allowbreak &-1\le z<2,\allowbreak \\0,\allowbreak &\text{其他},\allowbreak \end{cases}\)


（23）（I）证明略。（II）\(\displaystyle \dfrac{\displaystyle 2}{\displaystyle n(n-1)}\)。

\end{document}
